\section*{第 \ref{ch:g7:identifying-substances} 章——鉴别物质}

\begin{solution}{exo:g7:identifying-substances:1}
石灰水检验：将气体通入澄清的石灰水中，
若气体为二氧化碳，石灰水会变浑浊。
\end{solution}

\begin{solution}{exo:g7:identifying-substances:2}
粉末为无水硫酸铜，液体中含有水。
\end{solution}

\begin{solution}{exo:g7:identifying-substances:3}
氧气。
\end{solution}

\begin{solution}{exo:g7:identifying-substances:4}
铅笔线不溶于\omterm{def:g6:solutions-and-solubility:solute}{溶剂}。若用墨水画线，墨水线本身也会
被\omterm{def:g6:solutions-and-solubility:solute}{溶剂}带起并分离，破坏\omterm{def:g7:identifying-substances:chromatography}{色谱图}。
\end{solution}

\begin{solution}{exo:g7:identifying-substances:5}
若斑点位于\omterm{def:g6:solutions-and-solubility:solute}{溶剂}液面之下，染料会\omterm{def:g2:mixing-and-dissolving:dissolve}{溶解}到
烧杯中的\omterm{def:g6:solutions-and-solubility:solute}{溶剂}里，而不会沿滤纸向上移动。
\end{solution}

\begin{solution}{exo:g7:identifying-substances:6}
不是。它是两种染料的\omterm{def:g6:pure-substances-and-mixtures:mixture}{混合物}，
很可能是参考样品中的黄色染料和蓝色染料。
\end{solution}

\begin{solution}{exo:g7:identifying-substances:7}
用试管收集少量该气体，将点燃的木条放在
管口处：发出尖锐爆鸣声，说明为氢气。
\end{solution}

\begin{solution}{exo:g7:identifying-substances:8}
呼出的气体中二氧化碳的含量远高于新鲜\omterm{def:g4:air-a-mixture-of-gases:air}{空气}。
\end{solution}

\begin{solution}{exo:g7:identifying-substances:9}
该检验只能说明含有水，而不能说明只有水：盐水、
果汁或任何含水的\omterm{def:g6:pure-substances-and-mixtures:mixture}{混合物}同样会使粉末变蓝。
\end{solution}

\begin{solution}{exo:g7:identifying-substances:10}
三种染料（三个斑点）。黄色染料上升得最高。
\end{solution}

\begin{solution}{exo:g7:identifying-substances:11}
例如：将带火星的木条分别伸入各集气瓶，能使木条复燃的
是氧气。对其余两瓶，将点燃的木条置于瓶口，发出尖锐
爆鸣声的是氢气。剩下的一瓶是二氧化碳（可用石灰水确认）。
两三次检验即可。
\end{solution}

\begin{solution}{exo:g7:identifying-substances:12}
该气体可能不含二氧化碳，或含量太少，在
短时间内无法显现。对照实验用呼出的气体
（已知含有二氧化碳），可检验石灰水是否有效。
\end{solution}

\begin{solution}{pb:g7:identifying-substances:1}
\textbf{1.} 为使四者的条件（滤纸、\omterm{def:g6:solutions-and-solubility:solute}{溶剂}、时间）相同，
只有这样，才能比较斑点的高度。

\textbf{2.} \omterm{def:g6:pure-substances-and-mixtures:mixture}{混合物}：其墨水形成三个独立的斑点，即三种染料。

\textbf{3.} 两种染料。否。便条上的墨水含三种染料，
而 1 号笔既无红色染料，也无黄色染料。

\textbf{4.} 其三个斑点的高度与便条的不同，
说明它们是三种不同的染料。

\textbf{5.} 2 号笔：其三个斑点的高度
与便条上三个斑点的高度完全一致。

\textbf{6.} 氧气。

\textbf{7.} 氢气。

\textbf{8.} 二氧化碳。

\textbf{9.} 空白对照。它说明石灰水在这段时间内不会自行变浑浊，
通入普通\omterm{def:g4:air-a-mixture-of-gases:air}{空气}也不会变浑浊，因此
气体 Z 使石灰水变浑浊确实来自二氧化碳。

\textbf{10.} “水——无水硫酸铜——白色变蓝色”：
液体中含有水。

\textbf{11.} 不对。阳性检验只说明存在水，并不能说明
没有其他成分：该液体可能是溶液或
另一种含水\omterm{def:g6:pure-substances-and-mixtures:mixture}{混合物}。

\textbf{12.} 便条上的墨水含有 3 种染料，
且便条是由 2 号笔写成的。
\end{solution}
